\begin{itemize}
\item O computador envia 5 bytes: prescaler, frequência solicitada e quantidade de amostras.
\item O ATmega2560 responde com um byte de status e, em caso de sucesso, o bloco de dados.
\item UART0 em 8N1 e 115200 baud solicitados. A interface da placa faz a conexão USB com o computador.
\end{itemize}
\begin{block}{Custo da transmissão}
Cada amostra ocupa 1 byte, mas usa 10 bits no enlace:
\[
T_{envio}\approx\frac{4096\cdot10}{115200}=0{,}356\ \text{s}.
\]
\end{block}
\small Em \texttt{U2X0=1}, o divisor inteiro gera aproximadamente 117647 baud. A aquisição em bloco evita que a taxa serial limite diretamente $f_s$, mas há intervalo entre capturas.
\pblfonte{Firmware, UART0\_Init; relatório técnico, seção 3.3; esquema do Arduino Mega 2560.}
